% --- Abstract (max ~0.5 page, no tables/figures)
\newpage
\section*{Abstract}
\addcontentsline{toc}{section}{Abstract}

Space debris is a growing danger for satellites, and solar panels are the
parts that get hit most often because they are so large. In this project,
\textit{Space Debris Impact on Solar Panels}, we use conservation of linear
momentum, $\mathbf{P}_{\mathrm{before}} = \mathbf{P}_{\mathrm{after}}$, to see
how a small debris fragment changes a satellite's motion when it hits a
panel. We want to find the impulse the satellite gets, its change in
velocity, the spin it starts to have, and how much energy is lost, for
different impact conditions. Our hypothesis is that the velocity change is
tiny compared with the orbital speed, but the spin and the energy loss are
not. A Python program uses the impulse--momentum theorem and conservation of
linear and angular momentum to solve three scenarios, and it also runs
sweeps over debris mass, impact speed, impact angle and mass ratio. We model
the hit in two dimensions as an instant collision between a point-mass
fragment and a rigid satellite. Two ideal cases are solved: an elastic
collision with no friction, where the fragment bounces off, and a completely
inelastic collision, where the fragment sticks to the panel. The spin comes
from conservation of angular momentum about the centre of mass, and any
kinetic energy that is not left as motion counts as lost. In the baseline
case, a 1~g fragment at $10\unit{km\,s^{-1}}$ and $30^\circ$ sticks to the
panel 4~m from the centre of mass of a 500~kg satellite. The satellite's
velocity changes by only $20.0\unit{mm\,s^{-1}}$ ($2.6\times10^{-6}$ of its
orbital speed), but it starts spinning at $1.65^\circ\unit{s^{-1}}$, and
$99.999\%$ of the $50\unit{kJ}$ impact energy is lost as heat and damage.
Momentum, angular momentum and energy balance to about $10^{-16}$. The sweeps
show that the velocity change grows with debris mass and speed. For a
bouncing fragment it gets smaller as the hit becomes more grazing, while a
sticking fragment gives the same velocity change at every angle. Our model
only covers the momentum exchange right at the hit; holes, cracks and other
damage to the panel are not included.
